% ============================================================
% CHAPTER 9: FRONTEND SYSTEM ARCHITECTURE AND IMPLEMENTATION
% ============================================================


The AgriVerify frontend serves as the critical presentation layer, bridging deep learning inference engines and backend microservices. Engineered for high responsiveness, it delivers an intuitive interface tailored for diverse stakeholders—from farmers operating in low-resource environments to state administrators overseeing supply chains \cite{singh2026impact}. To meet enterprise scalability and real-time synchronization requirements, the architecture is constructed using Next.js~14, React~18, TypeScript~5, and Tailwind CSS. This stack enforces strict presentation logic separation, ensuring maintainability and robust cross-platform accessibility \cite{salim2024microservices}.

% ============================================================
\section{Frontend Architecture and Design Patterns}
\label{sec:frontend-architecture}

The presentation architecture adheres to layered design principles, isolating UI rendering from data synchronization, API communication, and authorization domains.

\subsection{Architectural Layers}
The request lifecycle processes through four specialized layers before interacting with backend services (Figure~\ref{fig:frontend_architecture}):
\begin{itemize}
    \item \textbf{Presentation UI (React \& Next.js):} Manages component rendering, form validation, and responsive layouts \cite{react2024}.
    \item \textbf{Client-Side State (TanStack Query):} Orchestrates asynchronous data fetching, intelligent caching, and optimistic mutations \cite{tanstack2024}.
    \item \textbf{Service Client (Axios):} Encapsulates network operations, standardizing payloads and error handling.
    \item \textbf{API Proxy (Next.js Routes):} Acts as a secure intermediary, routing client requests to ASP.NET Core endpoints while hiding internal URLs.
\end{itemize}

\begin{figure}[H]
\centering
\begin{adjustbox}{max width=0.95\textwidth}
\begin{tikzpicture}[
    node distance=0.6cm and 1.2cm,
    every node/.style={
        draw=black!80, thick, rounded corners=3pt, align=center,
        minimum width=4cm, minimum height=0.9cm, font=\scriptsize, fill=white
    },
    arrow/.style={->, thick, >=stealth, draw=blue!80!black}
]
\node[fill=blue!10] (ui) {\textbf{UI Component Layer}\\(Next.js \& React)};
\node[fill=green!10, below=of ui] (hooks) {\textbf{State Synchronization}\\(TanStack Query)};
\node[fill=yellow!10, below=of hooks] (services) {\textbf{Service Layer}\\(Axios API Client)};
\node[fill=orange!10, below=of services] (proxy) {\textbf{API Proxy Route}\\(Next.js Intermediary)};
\node[fill=red!10, right=1.2cm of services] (backend) {\textbf{ASP.NET Backend}\\(Business \& Data)};
\node[fill=purple!10, right=1.2cm of proxy] (ml) {\textbf{AI Inference}\\(Model Engine)};

\draw[arrow] (ui) -- (hooks);
\draw[arrow] (hooks) -- (services);
\draw[arrow] (services) -- (proxy);
\draw[arrow] (proxy) -| (backend);
\draw[arrow] (backend) -- (ml);
\end{tikzpicture}
\end{adjustbox}
\caption{Frontend Layered Architecture and Request Execution Flow}
\label{fig:frontend_architecture}
\end{figure}

\subsection{Secure Proxy Architecture}
To neutralize Cross-Origin Resource Sharing (CORS) vulnerabilities, outbound communications are routed through internal Next.js serverless endpoints (\texttt{/api/proxy/*}). This endpoint obfuscation conceals internal microservice topologies, keeps cryptographic secrets server-side, and normalizes security headers before browser delivery \cite{masue2024quantifying}.

% ============================================================
\section{Technology Stack Analysis}
\label{sec:technology-stack}

The technology stack is optimized for static type safety, developer velocity, and runtime execution speed (Table~\ref{tab:frontend-stack}).

\begin{table}[H]
\centering
\footnotesize
\caption{Frontend Technology Stack and Functional Rationale}
\label{tab:frontend-stack}
\begin{tabularx}{\textwidth}{@{}l X@{}}
\toprule
\textbf{Technology} & \textbf{Architectural Rationale} \\
\midrule
Next.js 14 & Enables Server-Side Rendering (SSR) and built-in API proxy routing \cite{nextjs2024}. \\
React 18 & Facilitates declarative UI rendering and efficient Virtual DOM reconciliation. \\
TypeScript 5 & Enforces strict compile-time type checking and backend DTO parity. \\
Tailwind CSS & Utility-first styling optimized for minimal production bundle sizes \cite{tailwind2024}. \\
Axios \& TanStack & Manages HTTP pipelines, token injection, and cache synchronization \cite{tanstack2024}. \\
Zustand 4.5 & Provides lightweight atomic state slices for session persistence. \\
\bottomrule
\end{tabularx}
\end{table}

% ============================================================
\section{Modular Project Structure}
\label{sec:project-structure}

To prevent architectural drift, the source tree utilizes strict domain-driven modularity. Key directories include \texttt{src/app} (Next.js routing), \texttt{src/components} (reusable UI primitives), \texttt{src/hooks} (data synchronization logic), and \texttt{src/types} (TypeScript DTO contracts matching backend entities). This decoupling isolates business logic from presentation constraints.

% ============================================================
\section{Routing Topology and Access Control}
\label{sec:routing-access}

The frontend leverages Role-Based Access Control (RBAC) to enforce security across public domains and protected administrative boundaries.

\subsection{Role Partitioning}
The application supports four hierarchical profiles: \textbf{Farmers} (verification access and dispute lodging), \textbf{Agricultural Officers} (inspection queues and evidence review), \textbf{Administrators} (system governance and telemetry), and \textbf{Public Users} (anonymous checks). 

\subsection{Route Protection Middleware}
A client-side \texttt{RoleGuard} component prevents privilege escalation by evaluating JSON Web Token (JWT) claims against required permissions before mounting React components.

\begin{lstlisting}[
    language=TypeScript,
    caption={Core Logic of the Role-Based Route Protection Guard},
    label={lst:role_guard},
    basicstyle=\ttfamily\scriptsize,
    breaklines=true,
    frame=single
]
export const RoleGuard: React.FC<{ children: React.ReactNode; allowedRoles: Role[] }> = ({ children, allowedRoles }) => {
    const { token, role, isLoading } = useAuthStore();
    const router = useRouter();

    useEffect(() => {
        if (!isLoading) {
            if (!token) return router.replace('/login');
            if (role && !allowedRoles.includes(role)) return router.replace('/unauthorized');
        }
    }, [token, role, isLoading, router, allowedRoles]);

    if (isLoading || !token) return <Loader />;
    return <>{children}</>;
};
\end{lstlisting}

% ============================================================
\section{Authentication and Secure Persistence}
\label{sec:authentication}

The platform secures user sessions using a dual-token JWT exchange, protecting against common client-side vulnerabilities like Cross-Site Scripting (XSS) \cite{masue2024quantifying}.

\subsection{Token Storage Architecture}
Short-lived access tokens are strictly held in the in-memory Zustand store to prevent malicious script extraction. Long-lived refresh tokens are persisted in encrypted caches (or \texttt{HttpOnly} cookies), facilitating background session renewal without compromising token integrity.

\subsection{Automated Token Rotation}
To maintain seamless UX during field operations, Axios interceptors autonomously detect \texttt{401 Unauthorized} responses, pause the network queue, execute a token refresh cycle, and seamlessly replay pending requests.

\begin{lstlisting}[
    language=TypeScript,
    caption={Axios Interceptor for Automated JWT Token Rotation},
    label={lst:axios_interceptor},
    basicstyle=\ttfamily\scriptsize,
    breaklines=true,
    frame=single
]
apiClient.interceptors.response.use(res => res, async (error) => {
    const origReq = error.config;
    if (error.response?.status === 401 && !origReq._retry) {
        origReq._retry = true;
        try {
            const res = await axios.post('/api/proxy/auth/refresh', { token: useAuthStore.getState().refreshToken });
            useAuthStore.getState().setTokens(res.data.accessToken, res.data.refreshToken);
            origReq.headers['Authorization'] = `Bearer ${res.data.accessToken}`;
            return apiClient(origReq);
        } catch (err) {
            useAuthStore.getState().logout();
            return Promise.reject(err);
        }
    }
    return Promise.reject(error);
});
\end{lstlisting}

% ============================================================
\section{API Integration and Type Safety}
\label{sec:api-integration}

To eliminate data corruption and runtime mapping errors, strict TypeScript Data Transfer Objects (DTOs) enforce compile-time verification of HTTP payloads. Integration modules (Auth, Verification, Analytics, and Complaint Services) mirror the ASP.NET Core backend schema, ensuring zero-defect data rendering across dashboards.

% ============================================================
\section{State Management}
\label{sec:state-management}

State architecture is bifurcated based on data volatility:

\subsection{Asynchronous Server State (TanStack Query)}
Network-driven data (e.g., dispute queues, telemetry) is managed by TanStack Query. It handles stale-while-revalidate caching, background polling, and optimistic UI updates, drastically reducing perceived network latency \cite{tanstack2024}.

\subsection{Synchronous Client State (Zustand)}
Client-local states, such as active UI themes and user session contexts, utilize Zustand. Its atomic state slices minimize React re-rendering cycles, conserving device battery and memory in resource-constrained field environments.

% ============================================================
\section{Responsive UI Design and Field Usability}
\label{sec:responsive-ui}

Recognizing the hardware limitations and environmental challenges of rural agriculture, the UI adopts a strict mobile-first design system utilizing Tailwind CSS \cite{singh2026impact}.

\subsection{Standardized UI Library}
To accelerate development and maintain cohesion, the system uses reusable primitives (Table~\ref{tab:ui-components}).

\begin{table}[H]
\centering
\footnotesize
\caption{Standardized UI Components}
\label{tab:ui-components}
\begin{tabularx}{\textwidth}{@{}l X@{}}
\toprule
\textbf{Component} & \textbf{Interaction Role} \\
\midrule
\textbf{Data Table} & Scrollable grids with sticky headers and dynamic pagination for inspection data. \\
\textbf{Status Badge} & Semantic color indicators (e.g., Green/Genuine) for rapid cognitive parsing. \\
\textbf{Loading Skeleton} & Animated structural placeholders preventing layout shifts during network latency. \\
\textbf{Input Modal} & Focused workflow overlays that trap keyboard navigation for accessible data entry. \\
\bottomrule
\end{tabularx}
\end{table}

\subsection{Accessibility Enhancements}
Usability optimizations include strict adherence to WCAG AAA contrast ratios for visibility under direct sunlight, enlarged touch targets ($48\times 48$px minimum) to minimize physical input errors, and full ARIA semantic tagging for assistive screen-reader compatibility.

% ============================================================
\section{Summary}
\label{sec:frontend-summary}

This chapter outlined the AgriVerify frontend architecture. By integrating Next.js, TypeScript, and state-of-the-art token rotation mechanisms, the platform delivers a secure, accessible, and performant user experience. Its decoupled layered design and mobile-first responsive framework successfully overcome the operational constraints of rural field deployments, bridging advanced AI analytics with practical agricultural workflows.